\section{Conclusiones}

    Se identificaron los requerimientos funcionales y no funcionales necesarios para el desarrollo de un aplicativo de realidad aumentada orientado al mantenimiento preventivo mediante la revisión de literatura, normas técnicas y el concepto brindado por personas con experiencia en el campo del mantenimiento. Este proceso permitió definir una solución alineada con las necesidades en contextos donde se trabaja bajo la metodología del mantenimiento preventivo, estableciendo criterios para el diseño de la interfaz, la gestión de la información técnica y la integración de experiencias de realidad aumentada.

    Se diseñó un aplicativo móvil que integra la consulta de información técnica, la gestión de documentos, uso de modelos tridimensionales y la ejecución de procedimientos de mantenimiento mediante realidad aumentada. Se desarrolló una estructura que permitió integrar estas funcionalidades en una plataforma única, manteniendo una organización por módulos facilitando su adaptación y futura ampliación. 

    La evaluación realizada sobre el compresor de aire de 3 HP permitió verificar el funcionamiento del aplicativo para un caso de visualización de una rutina de mantenimiento. Las pruebas confirmaron la correcta creación de marcadores, la calibración de modelos 3D, configuración de actividades y la reproducción de las guías visuales, evidenciando la ventaja de utilizar la realidad aumentada como herramienta de apoyo durante la capacitación y ejecución de tareas de mantenimiento preventivo al permitir la comprensión y seguimiento del proceso a realizar.

    La realización de pruebas en dispositivos móviles de diferentes gamas evidenció que el aplicativo puede ejecutarse correctamente en equipos compatibles con la tecnología de realidad aumentada. Aunque se observaron diferencias en los tiempos de inicio y navegación entre dispositivos, el comportamiento obtenido fue consistente con las capacidades de procesamiento de cada uno, permitiendo el funcionamiento de las principales características implementadas sin la generación de errores o cierres inesperados del aplicativo.

    Las herramientas que permiten configurar procedimientos de mantenimiento completos mediante la creación de marcadores, la calibración de modelos tridimensionales y la definición de secuencias de operaciones, convierten al aplicativo en una plataforma configurable para la generación de nuevas guías visuales, ampliando significativamente su potencial para diferentes tipos de equipos en diversos contextos industriales.